%% Font size %%
\documentclass[11pt]{article}

%% Load the custom package
\usepackage{Mathdoc}

%% Numéro de séquence %% Titre de la séquence %%
\renewcommand{\centerhead}{Chapitre 01 : Nombres relatifs - DS2 : Signe, distance à zéro, opposé}

%% Spacing commands %%
\renewcommand{\baselinestretch}{1} \setlength{\parindent}{0pt}

\begin{document}

\entetedevoirs{20}

\begin{exercicedevoir}[5]
Un nombre relatif est soit \textbf{positif}, soit \textbf{négatif} (0 étant à la fois positif
et négatif).\\
Indiquer, pour chaque nombre, s'il est \textbf{positif} ou \textbf{négatif}.
\smallskip

\begin{multicols}{2}
\begin{enumerate}[itemsep=0.6em]
\item $(+7)$ : \dotfill
\item $(-4)$ : \dotfill
\item $9$ : \dotfill
\item $(-12,5)$ : \dotfill
\item $(-0,3)$ : \dotfill
\item $0$ : \dotfill
\item $(+150)$ : \dotfill
\item $(-28)$ : \dotfill
\item $6,4$ : \dotfill
\item $(-1)$ : \dotfill
\end{enumerate}
\end{multicols}
\end{exercicedevoir}

\begin{exercicedevoir}[7]
Un nombre relatif est formé d'un \textbf{signe} ($+$ ou $-$) et d'une \textbf{distance à
zéro}.\\
Compléter le tableau suivant.
\smallskip

\begin{center}
\renewcommand{\arraystretch}{1.8}
\begin{tabular}{|c|c|c|c|c|c|c|c|}
\hline
Nombre & $(+6)$ & $(-9)$ & $(-14,2)$ & $(+3,5)$ & $(-100)$ & $8$ & $(-0,7)$ \\
\hline
Signe & & & & & & & \\
\hline
Distance à zéro & & & & & & & \\
\hline
\end{tabular}
\end{center}
\end{exercicedevoir}

\begin{exercicedevoir}[8]
Deux nombres relatifs sont \textbf{opposés} lorsqu'ils ont la même distance à zéro mais des
signes différents.\\
Donner l'opposé de chacun des nombres suivants.
\smallskip

\begin{multicols}{2}
\begin{enumerate}[itemsep=1em]
\item Opposé de $(+5)$ : \dotfill
\item Opposé de $(-9)$ : \dotfill
\item Opposé de $12$ : \dotfill
\item Opposé de $(-3,6)$ : \dotfill
\item Opposé de $(-27)$ : \dotfill
\item Opposé de $0,8$ : \dotfill
\item Opposé de $(-150)$ : \dotfill
\item Opposé de $41$ : \dotfill
\end{enumerate}
\end{multicols}
\end{exercicedevoir}

\end{document}
